\Chapter{59}

\Verse{1}
{
    \zA{话说宝玉闻听贾母等回来，随多添了一件衣裳，拄了杖前边来，都见过了。}
}
{
    \eA{[English source not present in the checked-in Bencraft/Joly files.]}
}
{
}

\Verse{2}
{
    \zA{贾 母等因每日辛苦，都要早些歇息，一宿无话。}
}
{
    \eA{[English source not present in the checked-in Bencraft/Joly files.]}
}
{
}

\Verse{3}
{
    \zA{次日五鼓，又往朝中去。}
}
{
    \eA{[English source not present in the checked-in Bencraft/Joly files.]}
}
{
}

\Verse{4}
{
    \zA{离送灵日不远，鸳鸯、琥珀、翡翠、玻璃四人都忙着打点贾母之物，玉钏、彩 云、彩霞皆打点王夫人之物，当面查点与跟随的管事媳妇们。}
}
{
    \eA{[English source not present in the checked-in Bencraft/Joly files.]}
}
{
}

\Verse{5}
{
    \zA{跟随的一共大小六个 丫鬟，十个老婆媳妇子，男人不算。}
}
{
    \eA{[English source not present in the checked-in Bencraft/Joly files.]}
}
{
}

\Verse{6}
{
    \zA{连日收拾驮轿器械。}
}
{
    \eA{[English source not present in the checked-in Bencraft/Joly files.]}
}
{
}

\Verse{7}
{
    \zA{鸳鸯和玉钏儿皆不随去， 只看屋子。}
}
{
    \eA{[English source not present in the checked-in Bencraft/Joly files.]}
}
{
}

\Verse{8}
{
    \zA{一面先几日预备帐幔铺陈之物，先有四五个媳妇并几个男子领出来，坐 了几辆车绕过去，先至下处，铺陈安插等候。}
}
{
    \eA{[English source not present in the checked-in Bencraft/Joly files.]}
}
{
}

\Verse{9}
{
    \zA{临日贾母带着贾蓉媳妇坐一乘驮轿， 王夫人在后，亦坐一乘驮轿，贾珍骑马率领众家丁围护。}
}
{
    \eA{[English source not present in the checked-in Bencraft/Joly files.]}
}
{
}

\Verse{10}
{
    \zA{又有几辆大车与婆子丫鬟 等坐，并放些随换的衣包等件。}
}
{
    \eA{[English source not present in the checked-in Bencraft/Joly files.]}
}
{
}

\Verse{11}
{
    \zA{是日薛姨妈尤氏率领诸人直送至大门外方回。}
}
{
    \eA{[English source not present in the checked-in Bencraft/Joly files.]}
}
{
}

\Verse{12}
{
    \zA{贾琏 恐路上不便，一面打发他父母起身，赶上了贾母王夫人驮轿，自己也随后带领家丁 押后跟来。}
}
{
    \eA{[English source not present in the checked-in Bencraft/Joly files.]}
}
{
}

\Verse{13}
{
    \zA{荣府内，赖大添派人丁上夜，将两处厅院都关了，一应出入人等皆走西边小角 门，日落时便命关了仪门，不放人出入。}
}
{
    \eA{[English source not present in the checked-in Bencraft/Joly files.]}
}
{
}

\Verse{14}
{
    \zA{园中前后东西角门亦皆关锁，只留王夫人 大房之后常系他姐妹出入之门，东边通薛姨妈的角门，这两门因在里院，不必关锁。}
}
{
    \eA{[English source not present in the checked-in Bencraft/Joly files.]}
}
{
}

\Verse{15}
{
    \zA{里面鸳鸯和玉钏儿也将上房关了，自领丫鬟婆子下房去歇。}
}
{
    \eA{[English source not present in the checked-in Bencraft/Joly files.]}
}
{
}

\Verse{16}
{
    \zA{每日林之孝家的带领十 来个老婆子上夜，穿堂内又添了许多小厮打更，已安插得十分妥当。}
}
{
    \eA{[English source not present in the checked-in Bencraft/Joly files.]}
}
{
}

\Verse{17}
{
    \zA{一日清晓，宝钗春困已醒，搴帷下榻，微觉轻寒。}
}
{
    \eA{[English source not present in the checked-in Bencraft/Joly files.]}
}
{
}

\Verse{18}
{
    \zA{及启户视之，见院中土润苔 青，原来五更时落了几点微雨。}
}
{
    \eA{[English source not present in the checked-in Bencraft/Joly files.]}
}
{
}

\Verse{19}
{
    \zA{于是唤起湘云等人来，一面梳洗。}
}
{
    \eA{[English source not present in the checked-in Bencraft/Joly files.]}
}
{
}

\Verse{20}
{
    \zA{湘云因说两腮作 痒，恐又犯了桃花癣，因问宝钗要些蔷薇硝擦。}
}
{
    \eA{[English source not present in the checked-in Bencraft/Joly files.]}
}
{
}

\Verse{21}
{
    \zA{宝钗道：“前日剩的都给了琴妹妹 了。”}
}
{
    \eA{[English source not present in the checked-in Bencraft/Joly files.]}
}
{
}

\Verse{22}
{
    \zA{因说：“颦儿配了许多，我正要要他些来，因今年竟没发痒就忘了。”}
}
{
    \eA{[English source not present in the checked-in Bencraft/Joly files.]}
}
{
}

\Verse{23}
{
    \zA{因命 莺儿去取些来。}
}
{
    \eA{[English source not present in the checked-in Bencraft/Joly files.]}
}
{
}

\Verse{24}
{
    \zA{莺儿应了才去时，蕊官便说：“我和你去，顺便瞧瞧藕官。”}
}
{
    \eA{[English source not present in the checked-in Bencraft/Joly files.]}
}
{
}

\Verse{25}
{
    \zA{说着 径同莺儿出了蘅芜院。}
}
{
    \eA{[English source not present in the checked-in Bencraft/Joly files.]}
}
{
}

\Verse{26}
{
    \zA{二人你言我语，一面行走一面说笑，不觉到了柳叶渚。}
}
{
    \eA{[English source not present in the checked-in Bencraft/Joly files.]}
}
{
}

\Verse{27}
{
    \zA{顺着柳堤走来，因见叶 才点碧，丝若垂金，莺儿便笑道：“你会拿这柳条子编东西不会？”}
}
{
    \eA{[English source not present in the checked-in Bencraft/Joly files.]}
}
{
}

\Verse{28}
{
    \zA{蕊官笑道：“编 什么东西？”}
}
{
    \eA{[English source not present in the checked-in Bencraft/Joly files.]}
}
{
}

\Verse{29}
{
    \zA{莺儿道：“什么编不得？}
}
{
    \eA{[English source not present in the checked-in Bencraft/Joly files.]}
}
{
}

\Verse{30}
{
    \zA{玩的使的都可。}
}
{
    \eA{[English source not present in the checked-in Bencraft/Joly files.]}
}
{
}

\Verse{31}
{
    \zA{等我摘些下来，带着这叶子 编一个花篮，掐了各色花儿放在里头，才是好玩呢。”}
}
{
    \eA{[English source not present in the checked-in Bencraft/Joly files.]}
}
{
}

\Verse{32}
{
    \zA{说着且不去取硝，只伸手采 了许多嫩条命蕊官拿着，他却一行走一行编花篮。}
}
{
    \eA{[English source not present in the checked-in Bencraft/Joly files.]}
}
{
}

\Verse{33}
{
    \zA{随路见花便采一二枝，编出一个 玲珑过梁的篮子。}
}
{
    \eA{[English source not present in the checked-in Bencraft/Joly files.]}
}
{
}

\Verse{34}
{
    \zA{枝上自有本来翠叶满布，将花放上，却也别致有趣。}
}
{
    \eA{[English source not present in the checked-in Bencraft/Joly files.]}
}
{
}

\Verse{35}
{
    \zA{喜得蕊官笑 说：“好姐姐，给了我罢。”}
}
{
    \eA{[English source not present in the checked-in Bencraft/Joly files.]}
}
{
}

\Verse{36}
{
    \zA{莺儿道：“这一个送咱们林姑娘，回来咱们再多采些， 编几个大家玩。”}
}
{
    \eA{[English source not present in the checked-in Bencraft/Joly files.]}
}
{
}

\Verse{37}
{
    \zA{说着来至潇湘馆中。}
}
{
    \eA{[English source not present in the checked-in Bencraft/Joly files.]}
}
{
}

\Verse{38}
{
    \zA{黛玉也正晨妆，见了这篮子，便笑说：“这 个新鲜花篮是谁编的？”}
}
{
    \eA{[English source not present in the checked-in Bencraft/Joly files.]}
}
{
}

\Verse{39}
{
    \zA{莺儿说：“我编的，送给姑娘玩的。”}
}
{
    \eA{[English source not present in the checked-in Bencraft/Joly files.]}
}
{
}

\Verse{40}
{
    \zA{黛玉接了，笑道： “怪道人人赞你的手巧，这玩意儿却也别致。”}
}
{
    \eA{[English source not present in the checked-in Bencraft/Joly files.]}
}
{
}

\Verse{41}
{
    \zA{一面瞧了，一面便叫紫鹃挂在那里。}
}
{
    \eA{[English source not present in the checked-in Bencraft/Joly files.]}
}
{
}

\Verse{42}
{
    \zA{莺儿又问候薛姨妈，方和黛玉要硝。}
}
{
    \eA{[English source not present in the checked-in Bencraft/Joly files.]}
}
{
}

\Verse{43}
{
    \zA{黛玉忙命紫鹃去包了一包，递给莺儿。}
}
{
    \eA{[English source not present in the checked-in Bencraft/Joly files.]}
}
{
}

\Verse{44}
{
    \zA{黛玉又 说道：“我好了，今日要出去逛逛。}
}
{
    \eA{[English source not present in the checked-in Bencraft/Joly files.]}
}
{
}

\Verse{45}
{
    \zA{你回去说给姐姐，不用过来问候妈妈，也不敢 劳他过来。}
}
{
    \eA{[English source not present in the checked-in Bencraft/Joly files.]}
}
{
}

\Verse{46}
{
    \zA{我梳了头，和妈妈都往那里去吃饭，大家热闹些。”}
}
{
    \eA{[English source not present in the checked-in Bencraft/Joly files.]}
}
{
}

\Verse{47}
{
    \zA{莺儿答应了出来，便到紫鹃房中找蕊官。}
}
{
    \eA{[English source not present in the checked-in Bencraft/Joly files.]}
}
{
}

\Verse{48}
{
    \zA{只见蕊官却与藕官二人正说得高兴， 不能相舍，莺儿便笑说：“姑娘也去呢，藕官先同去等着不好吗？”}
}
{
    \eA{[English source not present in the checked-in Bencraft/Joly files.]}
}
{
}

\Verse{49}
{
    \zA{紫鹃听见如此 说，便也说道：“这话倒很是。}
}
{
    \eA{[English source not present in the checked-in Bencraft/Joly files.]}
}
{
}

\Verse{50}
{
    \zA{他这里淘气的可厌。”}
}
{
    \eA{[English source not present in the checked-in Bencraft/Joly files.]}
}
{
}

\Verse{51}
{
    \zA{一面说，一面便将黛玉的匙 箸用了一块洋巾包了交给藕官，道：“你先带了这个去，也算一趟差了。”}
}
{
    \eA{[English source not present in the checked-in Bencraft/Joly files.]}
}
{
}

\Verse{52}
{
    \zA{藕官接 了，笑嘻嘻同他二人出来，一径顺着柳堤走来。}
}
{
    \eA{[English source not present in the checked-in Bencraft/Joly files.]}
}
{
}

\Verse{53}
{
    \zA{莺儿便又采些柳条，索性坐在山石 上编起来，又命蕊官先送了硝去再来。}
}
{
    \eA{[English source not present in the checked-in Bencraft/Joly files.]}
}
{
}

\Verse{54}
{
    \zA{他二人只顾爱看他编，那里舍得去？}
}
{
    \eA{[English source not present in the checked-in Bencraft/Joly files.]}
}
{
}

\Verse{55}
{
    \zA{莺儿只 管催，说：“你们再不去，我就不编了。”}
}
{
    \eA{[English source not present in the checked-in Bencraft/Joly files.]}
}
{
}

\Verse{56}
{
    \zA{藕官便说：“同你去了，再快回来。”}
}
{
    \eA{[English source not present in the checked-in Bencraft/Joly files.]}
}
{
}

\Verse{57}
{
    \zA{二人方去了。}
}
{
    \eA{[English source not present in the checked-in Bencraft/Joly files.]}
}
{
}

\Verse{58}
{
    \zA{这里莺儿正编，只见何妈的女儿春燕走来，笑问：“姐姐编什么呢？”}
}
{
    \eA{[English source not present in the checked-in Bencraft/Joly files.]}
}
{
}

\Verse{59}
{
    \zA{正说着， 蕊官藕官也到了，春燕便向藕官道：“前日你到底烧了什么纸？}
}
{
    \eA{[English source not present in the checked-in Bencraft/Joly files.]}
}
{
}

\Verse{60}
{
    \zA{叫我姨妈看见了， 要告你没告成，倒被宝玉赖了他好些不是，气得他一五一十告诉我妈。}
}
{
    \eA{[English source not present in the checked-in Bencraft/Joly files.]}
}
{
}

\Verse{61}
{
    \zA{你们在外头 二三年了，积了些什么仇恨，如今还不解开？”}
}
{
    \eA{[English source not present in the checked-in Bencraft/Joly files.]}
}
{
}

\Verse{62}
{
    \zA{藕官冷笑道：“有什么仇恨？}
}
{
    \eA{[English source not present in the checked-in Bencraft/Joly files.]}
}
{
}

\Verse{63}
{
    \zA{他们 不知足，反怨我们。}
}
{
    \eA{[English source not present in the checked-in Bencraft/Joly files.]}
}
{
}

\Verse{64}
{
    \zA{在外头这两年，不知赚了我们多少东西，你说说可有的没的？”}
}
{
    \eA{[English source not present in the checked-in Bencraft/Joly files.]}
}
{
}

\Verse{65}
{
    \zA{春燕也笑道：“他是我的姨妈，也不好向着外人反说他的。}
}
{
    \eA{[English source not present in the checked-in Bencraft/Joly files.]}
}
{
}

\Verse{66}
{
    \zA{怨不得宝玉说：‘女孩 儿未出嫁是颗无价宝珠，出了嫁不知怎么就变出许多不好的毛病儿来，再老了，更 不是珠子，竟是鱼眼睛了。}
}
{
    \eA{[English source not present in the checked-in Bencraft/Joly files.]}
}
{
}

\Verse{67}
{
    \zA{分明一个人，怎么变出三样来。’}
}
{
    \eA{[English source not present in the checked-in Bencraft/Joly files.]}
}
{
}

\Verse{68}
{
    \zA{这话虽是混帐话，想 起来真不错。}
}
{
    \eA{[English source not present in the checked-in Bencraft/Joly files.]}
}
{
}

\Verse{69}
{
    \zA{别人不知道，只说我妈和姨妈他老姐儿两个，如今越老了越把钱看的 真了。}
}
{
    \eA{[English source not present in the checked-in Bencraft/Joly files.]}
}
{
}

\Verse{70}
{
    \zA{先是老姐儿两个在家抱怨没个差使进益，幸亏有了这园子，把我挑进来。}
}
{
    \eA{[English source not present in the checked-in Bencraft/Joly files.]}
}
{
}

\Verse{71}
{
    \zA{可 巧把我分到怡红院，家里省了我一个人的费用不算外，每月还有四五百钱的馀剩， 这也还说不够。}
}
{
    \eA{[English source not present in the checked-in Bencraft/Joly files.]}
}
{
}

\Verse{72}
{
    \zA{后来老姐儿两个都派到梨香院去照看他们，藕官认了我姨妈，芳官 认了我妈，这几年着实宽绰了。}
}
{
    \eA{[English source not present in the checked-in Bencraft/Joly files.]}
}
{
}

\Verse{73}
{
    \zA{如今挪进来，也算撂开手了，还只无厌，你说可笑 不可笑？}
}
{
    \eA{[English source not present in the checked-in Bencraft/Joly files.]}
}
{
}

\Verse{74}
{
    \zA{接着我妈和芳官又吵了一场，又要给宝玉吹汤，讨个没趣儿。}
}
{
    \eA{[English source not present in the checked-in Bencraft/Joly files.]}
}
{
}

\Verse{75}
{
    \zA{幸亏园里的 人多，没人记的清楚谁是谁的亲故，要有人记得，我们一家子叫人家看着什么意思 呢。}
}
{
    \eA{[English source not present in the checked-in Bencraft/Joly files.]}
}
{
}

\Verse{76}
{
    \zA{你这会子又跑了来弄这个，这一带地方上的东西都是我姑妈管着。}
}
{
    \eA{[English source not present in the checked-in Bencraft/Joly files.]}
}
{
}

\Verse{77}
{
    \zA{他一得了这 地，每日起早睡晚自己辛苦了还不算，每日逼着我们来照看，生怕有人遭塌，―― 我又怕误了我的差使。}
}
{
    \eA{[English source not present in the checked-in Bencraft/Joly files.]}
}
{
}

\Verse{78}
{
    \zA{如今我们进来了，老姑嫂两个照看得谨谨慎慎，一根草也不 许人乱动。}
}
{
    \eA{[English source not present in the checked-in Bencraft/Joly files.]}
}
{
}

\Verse{79}
{
    \zA{你还掐这些好花儿，又折他的嫩树枝子，他们即刻就来，你看他们抱怨。”}
}
{
    \eA{[English source not present in the checked-in Bencraft/Joly files.]}
}
{
}

\Verse{80}
{
    \zA{莺儿道：“别人折掐使不得，独我使得。}
}
{
    \eA{[English source not present in the checked-in Bencraft/Joly files.]}
}
{
}

\Verse{81}
{
    \zA{自从分了地基之后，各房里每日皆有分例 的不用算，单算花草玩意儿：谁管什么，每日谁就把各房里姑娘丫头戴的，必要各 色送些折枝去，另有插瓶的。}
}
{
    \eA{[English source not present in the checked-in Bencraft/Joly files.]}
}
{
}

\Verse{82}
{
    \zA{惟有我们姑娘说了：‘一概不用送，等要什么再和你 要。’}
}
{
    \eA{[English source not present in the checked-in Bencraft/Joly files.]}
}
{
}

\Verse{83}
{
    \zA{究竟总没要过一次。}
}
{
    \eA{[English source not present in the checked-in Bencraft/Joly files.]}
}
{
}

\Verse{84}
{
    \zA{我今便掐些，他们也不好意思说的。”}
}
{
    \eA{[English source not present in the checked-in Bencraft/Joly files.]}
}
{
}

\Verse{85}
{
    \zA{一言未了，他姑妈果然拄了拐杖走来，莺儿春燕等忙让坐。}
}
{
    \eA{[English source not present in the checked-in Bencraft/Joly files.]}
}
{
}

\Verse{86}
{
    \zA{那婆子见采了许多 嫩柳，又见藕官等采了许多鲜花，心里便不受用，看着莺儿编弄，又不好说什么。}
}
{
    \eA{[English source not present in the checked-in Bencraft/Joly files.]}
}
{
}

\Verse{87}
{
    \zA{便说春燕道：“我叫你来照看照看，你就贪着玩不去了。}
}
{
    \eA{[English source not present in the checked-in Bencraft/Joly files.]}
}
{
}

\Verse{88}
{
    \zA{倘或叫起你来，你又说我 使你了，拿我作隐身草儿，你来乐！”}
}
{
    \eA{[English source not present in the checked-in Bencraft/Joly files.]}
}
{
}

\Verse{89}
{
    \zA{春燕道：“你老人家又使我，又怕，这会子 反说我，难道把我劈八瓣子不成？”}
}
{
    \eA{[English source not present in the checked-in Bencraft/Joly files.]}
}
{
}

\Verse{90}
{
    \zA{莺儿笑道：“姑妈，你别信小燕儿的话。}
}
{
    \eA{[English source not present in the checked-in Bencraft/Joly files.]}
}
{
}

\Verse{91}
{
    \zA{这都 是他摘下来，烦我给他编，我撵他，他不去。”}
}
{
    \eA{[English source not present in the checked-in Bencraft/Joly files.]}
}
{
}

\Verse{92}
{
    \zA{春燕笑道：“你可少玩儿!你只顾 玩，他老人家就认真的。”}
}
{
    \eA{[English source not present in the checked-in Bencraft/Joly files.]}
}
{
}

\Verse{93}
{
    \zA{那婆子本是愚夯之辈，兼之年迈昏？}
}
{
    \eA{[English source not present in the checked-in Bencraft/Joly files.]}
}
{
}

\Verse{94}
{
    \zA{，惟利是命，一概 情面不管。}
}
{
    \eA{[English source not present in the checked-in Bencraft/Joly files.]}
}
{
}

\Verse{95}
{
    \zA{正心疼肝断，无计可施，听莺儿如此说，便倚老卖老，拿起拄杖向春燕 身上击了几下，骂道：“小蹄子!我说着你，你还和我强嘴儿呢。}
}
{
    \eA{[English source not present in the checked-in Bencraft/Joly files.]}
}
{
}

\Verse{96}
{
    \zA{你妈恨的牙痒痒， 要撕你的肉吃呢，你还和我梆子似的！”}
}
{
    \eA{[English source not present in the checked-in Bencraft/Joly files.]}
}
{
}

\Verse{97}
{
    \zA{打得春燕又愧又急，因哭道：“莺儿姐姐 玩话，你就认真打我!我妈为什么恨我？}
}
{
    \eA{[English source not present in the checked-in Bencraft/Joly files.]}
}
{
}

\Verse{98}
{
    \zA{又没烧糊了洗脸水，有什么不是？”}
}
{
    \eA{[English source not present in the checked-in Bencraft/Joly files.]}
}
{
}

\Verse{99}
{
    \zA{莺儿本 是玩话，忽见婆子认真动了气，忙上前拉住，笑道：“我才是玩话，你老人家打他， 这不是臊我了吗？”}
}
{
    \eA{[English source not present in the checked-in Bencraft/Joly files.]}
}
{
}

\Verse{100}
{
    \zA{那婆子道：“姑娘你别管我们的事。}
}
{
    \eA{[English source not present in the checked-in Bencraft/Joly files.]}
}
{
}

\Verse{101}
{
    \zA{难道为姑娘在这里，不许 我们管孩子不成？”}
}
{
    \eA{[English source not present in the checked-in Bencraft/Joly files.]}
}
{
}

\Verse{102}
{
    \zA{莺儿听这般蠢话，便赌气红了脸，撒了手，冷笑道：“你要管， 那一刻管不得？}
}
{
    \eA{[English source not present in the checked-in Bencraft/Joly files.]}
}
{
}

\Verse{103}
{
    \zA{偏我说了一句玩话，就管他了？}
}
{
    \eA{[English source not present in the checked-in Bencraft/Joly files.]}
}
{
}

\Verse{104}
{
    \zA{我看你管去！”}
}
{
    \eA{[English source not present in the checked-in Bencraft/Joly files.]}
}
{
}

\Verse{105}
{
    \zA{说着便坐下，仍编柳 篮子。}
}
{
    \eA{[English source not present in the checked-in Bencraft/Joly files.]}
}
{
}

\Verse{106}
{
    \zA{偏又春燕的娘出来找他，喊道：“你不来舀水，在那里做什么？”}
}
{
    \eA{[English source not present in the checked-in Bencraft/Joly files.]}
}
{
}

\Verse{107}
{
    \zA{那婆子便接 声儿道：“你来瞧瞧!你女孩儿连我也不服了，在这里排揎我呢。”}
}
{
    \eA{[English source not present in the checked-in Bencraft/Joly files.]}
}
{
}

\Verse{108}
{
    \zA{那婆子一面走 过来，说：“姑奶奶又怎么了？}
}
{
    \eA{[English source not present in the checked-in Bencraft/Joly files.]}
}
{
}

\Verse{109}
{
    \zA{我们丫头眼里没娘罢了，连姑妈也没了不成？”}
}
{
    \eA{[English source not present in the checked-in Bencraft/Joly files.]}
}
{
}

\Verse{110}
{
    \zA{莺 儿见他娘来了，只得又说原故。}
}
{
    \eA{[English source not present in the checked-in Bencraft/Joly files.]}
}
{
}

\Verse{111}
{
    \zA{他姑娘那里容人说话？}
}
{
    \eA{[English source not present in the checked-in Bencraft/Joly files.]}
}
{
}

\Verse{112}
{
    \zA{便将石上的花柳与他娘瞧， 道：“你瞧瞧，你女孩儿这么大孩子顽的。}
}
{
    \eA{[English source not present in the checked-in Bencraft/Joly files.]}
}
{
}

\Verse{113}
{
    \zA{他领着人遭塌我，我怎么说人？”}
}
{
    \eA{[English source not present in the checked-in Bencraft/Joly files.]}
}
{
}

\Verse{114}
{
    \zA{他娘 也正为芳官之气未平，又恨春燕不遂他的心，便走上来打了个耳刮子，骂道：“小 娼妇，你能上了几年台盘，你也跟着那起轻薄浪小妇学!怎么就管不得你们了？}
}
{
    \eA{[English source not present in the checked-in Bencraft/Joly files.]}
}
{
}

\Verse{115}
{
    \zA{干的 我管不得，你是我自己生出来的，难道也不敢管你不成？}
}
{
    \eA{[English source not present in the checked-in Bencraft/Joly files.]}
}
{
}

\Verse{116}
{
    \zA{既是你们这起蹄子到得去 的地方我到不去，你就死在那里伺候，又跑出来浪汉子！”}
}
{
    \eA{[English source not present in the checked-in Bencraft/Joly files.]}
}
{
}

\Verse{117}
{
    \zA{一面又抓起那柳条子来， 直送到他脸上，问道：“这叫做什么？}
}
{
    \eA{[English source not present in the checked-in Bencraft/Joly files.]}
}
{
}

\Verse{118}
{
    \zA{这编的是你娘的什么？”}
}
{
    \eA{[English source not present in the checked-in Bencraft/Joly files.]}
}
{
}

\Verse{119}
{
    \zA{莺儿忙道：“那是 我编的，你别指桑骂槐的。”}
}
{
    \eA{[English source not present in the checked-in Bencraft/Joly files.]}
}
{
}

\Verse{120}
{
    \zA{那婆子深妒袭人晴雯一干人，早知道凡房中大些的丫 鬟，都比他们有些体统权势。}
}
{
    \eA{[English source not present in the checked-in Bencraft/Joly files.]}
}
{
}

\Verse{121}
{
    \zA{凡见了这一干人，心中又畏又让，未免又气又恨，亦 且迁怒于众；复又看见了藕官，又是他姐姐的冤家：四处凑成一股怒气。}
}
{
    \eA{[English source not present in the checked-in Bencraft/Joly files.]}
}
{
}

\Verse{122}
{
    \zA{那春燕啼哭着往怡红院去了。}
}
{
    \eA{[English source not present in the checked-in Bencraft/Joly files.]}
}
{
}

\Verse{123}
{
    \zA{他娘又恐问他为何哭，怕他又说出来，又要受晴 雯等的气，不免赶着来喊道：“你回来!我告诉你再去。”}
}
{
    \eA{[English source not present in the checked-in Bencraft/Joly files.]}
}
{
}

\Verse{124}
{
    \zA{春燕那里肯回来。}
}
{
    \eA{[English source not present in the checked-in Bencraft/Joly files.]}
}
{
}

\Verse{125}
{
    \zA{急的 他娘跑了去要拉他，春燕回头看见，便也往前飞跑。}
}
{
    \eA{[English source not present in the checked-in Bencraft/Joly files.]}
}
{
}

\Verse{126}
{
    \zA{他娘只顾赶他，不防脚下被青 苔滑倒。}
}
{
    \eA{[English source not present in the checked-in Bencraft/Joly files.]}
}
{
}

\Verse{127}
{
    \zA{招的莺儿三个人反都笑了。}
}
{
    \eA{[English source not present in the checked-in Bencraft/Joly files.]}
}
{
}

\Verse{128}
{
    \zA{莺儿赌气将花柳皆掷于河中，自回房去。}
}
{
    \eA{[English source not present in the checked-in Bencraft/Joly files.]}
}
{
}

\Verse{129}
{
    \zA{这里 把个婆子心疼的只念佛，又骂：“促狭小蹄子!遭塌了花儿，雷也是要劈的。”}
}
{
    \eA{[English source not present in the checked-in Bencraft/Joly files.]}
}
{
}

\Verse{130}
{
    \zA{自 己且掐花与各房送去。}
}
{
    \eA{[English source not present in the checked-in Bencraft/Joly files.]}
}
{
}

\Verse{131}
{
    \zA{却说春燕一直跑进院中，顶头遇见袭人往黛玉处问安去，春燕便一把抱住袭人 说：“姑娘救我，我妈又打我呢！”}
}
{
    \eA{[English source not present in the checked-in Bencraft/Joly files.]}
}
{
}

\Verse{132}
{
    \zA{袭人见他娘来了，不免生气，便说道：“三日 两头儿，打了干的打亲的。}
}
{
    \eA{[English source not present in the checked-in Bencraft/Joly files.]}
}
{
}

\Verse{133}
{
    \zA{还是卖弄你女孩儿多，还是认真不知王法？”}
}
{
    \eA{[English source not present in the checked-in Bencraft/Joly files.]}
}
{
}

\Verse{134}
{
    \zA{这婆子来 了几日，见袭人不言不语，是好性儿的，便说道：“姑娘，你不知道，别管我们的 闲事。}
}
{
    \eA{[English source not present in the checked-in Bencraft/Joly files.]}
}
{
}

\Verse{135}
{
    \zA{都是你们纵的，还管什么？”}
}
{
    \eA{[English source not present in the checked-in Bencraft/Joly files.]}
}
{
}

\Verse{136}
{
    \zA{说着，便又赶着打。}
}
{
    \eA{[English source not present in the checked-in Bencraft/Joly files.]}
}
{
}

\Verse{137}
{
    \zA{袭人气的转身进来，见麝 月正在海棠下晾手巾，――听如此喊闹，便说：“姐姐别管，看他怎么着。”}
}
{
    \eA{[English source not present in the checked-in Bencraft/Joly files.]}
}
{
}

\Verse{138}
{
    \zA{一面 使眼色给春燕。}
}
{
    \eA{[English source not present in the checked-in Bencraft/Joly files.]}
}
{
}

\Verse{139}
{
    \zA{春燕会意，直奔了宝玉去。}
}
{
    \eA{[English source not present in the checked-in Bencraft/Joly files.]}
}
{
}

\Verse{140}
{
    \zA{众人都笑说：“这可是!从来没有的事， 今儿都闹出来了。”}
}
{
    \eA{[English source not present in the checked-in Bencraft/Joly files.]}
}
{
}

\Verse{141}
{
    \zA{麝月向婆子道：“你再略煞一煞气儿，难道这些人的脸面，和 你讨一个情还讨不出来不成？”}
}
{
    \eA{[English source not present in the checked-in Bencraft/Joly files.]}
}
{
}

\Verse{142}
{
    \zA{那婆子见他女儿奔到宝玉身边去，又见宝玉拉了春燕的手，说：“你别怕，有 我呢。”}
}
{
    \eA{[English source not present in the checked-in Bencraft/Joly files.]}
}
{
}

\Verse{143}
{
    \zA{春燕一行哭，一行将方才莺儿等事都说出来。}
}
{
    \eA{[English source not present in the checked-in Bencraft/Joly files.]}
}
{
}

\Verse{144}
{
    \zA{宝玉越发急起来，说：“你 只在这里闹倒罢了，怎么把你妈也都得罪起来？”}
}
{
    \eA{[English source not present in the checked-in Bencraft/Joly files.]}
}
{
}

\Verse{145}
{
    \zA{麝月又向婆子及众人道：“怨不 得这嫂子说我们管不着他们的事。}
}
{
    \eA{[English source not present in the checked-in Bencraft/Joly files.]}
}
{
}

\Verse{146}
{
    \zA{我们原无知，错管了，如今请出一个管得着的人 来管一管，嫂子就心服口服，也知道规矩了。”}
}
{
    \eA{[English source not present in the checked-in Bencraft/Joly files.]}
}
{
}

\Verse{147}
{
    \zA{便回头命小丫头子：“去把平儿给 我叫来，平儿不得闲，就把林大娘叫了来。”}
}
{
    \eA{[English source not present in the checked-in Bencraft/Joly files.]}
}
{
}

\Verse{148}
{
    \zA{那小丫头子应了便走。}
}
{
    \eA{[English source not present in the checked-in Bencraft/Joly files.]}
}
{
}

\Verse{149}
{
    \zA{众媳妇上来笑 说：“嫂子快求姑娘们叫回那孩子来罢。}
}
{
    \eA{[English source not present in the checked-in Bencraft/Joly files.]}
}
{
}

\Verse{150}
{
    \zA{平姑娘来了，可就不好了。”}
}
{
    \eA{[English source not present in the checked-in Bencraft/Joly files.]}
}
{
}

\Verse{151}
{
    \zA{那婆子说道： “凭是那个姑娘来了，也要评个理。}
}
{
    \eA{[English source not present in the checked-in Bencraft/Joly files.]}
}
{
}

\Verse{152}
{
    \zA{没有见个娘管女孩儿，大家管着娘的！”}
}
{
    \eA{[English source not present in the checked-in Bencraft/Joly files.]}
}
{
}

\Verse{153}
{
    \zA{众人 笑道：“你当是那个平姑娘？}
}
{
    \eA{[English source not present in the checked-in Bencraft/Joly files.]}
}
{
}

\Verse{154}
{
    \zA{是二奶奶屋里的平姑娘啊。}
}
{
    \eA{[English source not present in the checked-in Bencraft/Joly files.]}
}
{
}

\Verse{155}
{
    \zA{他有情么，说你两句；他 一翻脸，嫂子你吃不了兜着走。”}
}
{
    \eA{[English source not present in the checked-in Bencraft/Joly files.]}
}
{
}

\Verse{156}
{
    \zA{说着只见那个小丫头回来说：“平姑娘正有事呢， 问我做什么，我告诉了他。}
}
{
    \eA{[English source not present in the checked-in Bencraft/Joly files.]}
}
{
}

\Verse{157}
{
    \zA{他说，叫先撵出他去，告诉林大娘，在角门子上打四十 板子就是了。”}
}
{
    \eA{[English source not present in the checked-in Bencraft/Joly files.]}
}
{
}

\Verse{158}
{
    \zA{那婆子听见如此说了，吓得泪流满面，央告袭人等说：“好容易我 进来了，况且我是寡妇家，没有坏心，一心在里头伏侍姑娘们。}
}
{
    \eA{[English source not present in the checked-in Bencraft/Joly files.]}
}
{
}

\Verse{159}
{
    \zA{我这一去，不知苦 到什么田地！”}
}
{
    \eA{[English source not present in the checked-in Bencraft/Joly files.]}
}
{
}

\Verse{160}
{
    \zA{袭人见他如此说，又心软了，便说：“你既要在这里，又不守规矩， 又不听话，又乱打人。}
}
{
    \eA{[English source not present in the checked-in Bencraft/Joly files.]}
}
{
}

\Verse{161}
{
    \zA{那里弄你这个不晓事的人来!天天斗口齿，也叫人笑话。”}
}
{
    \eA{[English source not present in the checked-in Bencraft/Joly files.]}
}
{
}

\Verse{162}
{
    \zA{晴雯道：“理他呢，打发他去了正经。}
}
{
    \eA{[English source not present in the checked-in Bencraft/Joly files.]}
}
{
}

\Verse{163}
{
    \zA{那里那么大工夫和他对嘴对舌的？”}
}
{
    \eA{[English source not present in the checked-in Bencraft/Joly files.]}
}
{
}

\Verse{164}
{
    \zA{那婆子 又央众人道：“我虽错了，姑娘们吩咐了，以后改过。}
}
{
    \eA{[English source not present in the checked-in Bencraft/Joly files.]}
}
{
}

\Verse{165}
{
    \zA{姑娘们那不是行好积德？”}
}
{
    \eA{[English source not present in the checked-in Bencraft/Joly files.]}
}
{
}

\Verse{166}
{
    \zA{一面又央告春燕：“原是为打你起的，饶没打成你，我如今反受了罪。}
}
{
    \eA{[English source not present in the checked-in Bencraft/Joly files.]}
}
{
}

\Verse{167}
{
    \zA{好孩子，你 好歹替我求求罢！”}
}
{
    \eA{[English source not present in the checked-in Bencraft/Joly files.]}
}
{
}

\Verse{168}
{
    \zA{宝玉见如此可怜，便命留下：“不许再闹!再闹，一定打了撵 出去。”}
}
{
    \eA{[English source not present in the checked-in Bencraft/Joly files.]}
}
{
}

\Verse{169}
{
    \zA{那婆子一一谢过下去。}
}
{
    \eA{[English source not present in the checked-in Bencraft/Joly files.]}
}
{
}

\Verse{170}
{
    \zA{只见平儿走来，问系何事，袭人等忙说：“已完了，不 必再提了。”}
}
{
    \eA{[English source not present in the checked-in Bencraft/Joly files.]}
}
{
}

\Verse{171}
{
    \zA{平儿笑道：“‘得饶人处且饶人’，得将就的就省些事罢。}
}
{
    \eA{[English source not present in the checked-in Bencraft/Joly files.]}
}
{
}

\Verse{172}
{
    \zA{但只听见 各屋里大小人等都作起反来了，一处不了又一处，叫我不知管那一处是。”}
}
{
    \eA{[English source not present in the checked-in Bencraft/Joly files.]}
}
{
}

\Verse{173}
{
    \zA{袭人笑 道：“我只说我们这里反了，原来还有几处。”}
}
{
    \eA{[English source not present in the checked-in Bencraft/Joly files.]}
}
{
}

\Verse{174}
{
    \zA{平儿笑道：“这算什么事!这三四 日的工夫，一共大小出了八九件呢，比这里的还大，可气可笑。”}
}
{
    \eA{[English source not present in the checked-in Bencraft/Joly files.]}
}
{
}

\Verse{175}
{
    \zA{袭人等听了诧异。}
}
{
    \eA{[English source not present in the checked-in Bencraft/Joly files.]}
}
{
}

\Verse{176}
{
    \zA{不知何事，下回分解。}
}
{
    \eA{[English source not present in the checked-in Bencraft/Joly files.]}
}
{
}

\Verse{177}
{
    \zA{(本章完)}
}
{
    \eA{[English source not present in the checked-in Bencraft/Joly files.]}
}
{
}
